\subsection{The Identification Problem}

Consider an encoder-decoder pair communicating about entities from a large universe $\mathcal{V}$. The decoder must \emph{identify} each entity, determining which of $k$ classes it belongs to, using only:
\begin{itemize}
\item A \emph{tag} of $L$ bits stored with the entity, and/or
\item \emph{Queries} to a binary oracle: ``does entity $v$ satisfy attribute $I$?''
\end{itemize}

This is not reconstruction (the decoder need not recover $v$), but \emph{identification} in the sense of Ahlswede and Dueck \cite{ahlswede1989identification}: the decoder must answer ``which class?'' with zero or bounded error. Our work extends this framework to consider the tradeoff between tag storage, query complexity, and identification accuracy.

We prove three core theorem families, and then extend them with finite-block, open-world, and mechanization results:
\begin{enumerate}
\item \textbf{Information barrier (identifiability limit).} When the attribute profile $\pi: \mathcal{V} \to \{0,1\}^n$ is not injective on classes, zero-error identification via queries alone is impossible: any decoder produces identical output on colliding classes, so cannot be correct for both.
\item \textbf{Optimal tagging (achievability).} A tag of $L = \lceil \log_2 k \rceil$ bits achieves zero-error identification with $W = O(1)$ query cost. For maximal-barrier domains ($A_\pi = k$), this is the unique Pareto-optimal point in the $(L, W, D)$ tradeoff space at $D=0$; in general domains, the converse depends on $A_\pi := \max_u |\{c : \pi(c)=u\}|$.
\item \textbf{Matroid structure (query complexity).} Minimal sufficient query sets form the bases of a matroid. The \emph{distinguishing dimension} (the common cardinality of all minimal sets) is well-defined and lower-bounds the query cost $W$ for any tag-free scheme.
\end{enumerate}

These results are domain-independent within the stated observation model: the same lower bounds appear across databases, biological taxonomy, knowledge graphs, and software runtimes. Once classes collide under observable attributes, constant-cost zero-error identification requires auxiliary naming information. In maximal-barrier domains this yields a unique $D=0$ Pareto solution, so the issue is resource feasibility, not domain-specific style conventions.

\subsection{The Observation Model}

We formalize the observational constraint as a family of binary predicates. The terminology is deliberately abstract; concrete instantiations follow in Section~\ref{sec:applications}.

\begin{definition}[Entity space and attribute family]
Let $\mathcal{V}$ be a set of entities (program objects, database records, biological specimens, library items). Let $\mathcal{I}$ be a finite set of binary \emph{attributes}. Each $I \in \mathcal{I}$ induces a bipartition of $\mathcal{V}$ via $q_I$, and the full family induces the observational equivalence partition.
\end{definition}

\begin{remark}[Terminology]
We use ``attribute'' for the abstract concept. Depending on domain, attributes may be interfaces/signatures, database columns, phenotypic characters, or catalog facets. The mathematics is identical.
\end{remark}

\begin{definition}[Attribute observation family]
For each $I \in \mathcal{I}$, define the attribute-membership observation $q_I: \mathcal{V} \to \{0,1\}$:
\[
q_I(v) = \begin{cases} 1 & \text{if } v \text{ satisfies attribute } I \\ 0 & \text{otherwise} \end{cases}
\]
Let $\Phi_{\mathcal{I}} = \{q_I : I \in \mathcal{I}\}$ denote the attribute observation family.
\end{definition}

\begin{remark}[Notation for size parameters]
We write $n := |\mathcal{I}|$ for the ambient number of available attributes. We write $d$ for the distinguishing dimension (the common size of all minimal distinguishing query sets; Definition~\ref{def:distinguishing-dimension}), so $d \le n$ and there exist worst-case families with $d = n$. We write $m$ for the number of \emph{query sites} (call sites) that perform attribute checks in a program or protocol (used only in the complexity-of-maintenance discussion).
When discussing a particular identification/verification task, we may write $s$ for the number of attributes actually queried/traversed by the procedure (e.g., interface or schema checks in software/data systems, phenotypic checks in taxonomy), with $s \le n$. The maintenance-only parameter $m$ appears only in Section~\ref{sec:complexity}.
\end{remark}

\begin{definition}[Attribute profile]
The attribute profile function $\pi: \mathcal{V} \to \{0,1\}^{|\mathcal{I}|}$ maps each value to its complete attribute signature:
\[
\pi(v) = (q_I(v))_{I \in \mathcal{I}}
\]
\end{definition}

\begin{definition}[Attribute indistinguishability]
Values $v, w \in \mathcal{V}$ are \emph{attribute-indistinguishable}, written $v \sim w$, iff $\pi(v) = \pi(w)$.
\end{definition}

The relation $\sim$ is an equivalence relation. We write $[v]_\sim$ for the equivalence class of $v$.

\begin{definition}[Attribute-only observer]
An \emph{attribute-only observer} is any procedure whose interaction with a value $v \in \mathcal{V}$ is limited to queries in $\Phi_{\mathcal{I}}$. Formally, the observer interacts with $v$ only via primitive attribute queries $q_I \in \Phi_{\mathcal{I}}$; hence any transcript (and output) factors through $\pi(v)$.
\end{definition}

\subsection{The Central Question}

The central question is: \textbf{what semantic properties can an attribute-only observer compute?}

A semantic property is a function $P: \mathcal{V} \to \{0,1\}$ (or more generally, $P: \mathcal{V} \to Y$ for some codomain $Y$). We say $P$ is \emph{attribute-computable} if there exists a function $f: \{0,1\}^{|\mathcal{I}|} \to Y$ such that $P(v) = f(\pi(v))$ for all $v$.

\subsection{The Information Barrier}

\begin{theorem}[Information barrier]\label{thm:information-barrier}
Let $P: \mathcal{V} \to Y$ be any function. If $P$ is attribute-computable, then $P$ is constant on $\sim$-equivalence classes:
\[
v \sim w \implies P(v) = P(w)
\]
Equivalently: no attribute-only observer can compute any property that varies within an equivalence class.
\end{theorem}
\leanmeta{ \LH{ACS8}, \LH{ACS9}, \LH{COH1}}

\begin{proof}
Suppose $P$ is attribute-computable via $f$, i.e., $P(v) = f(\pi(v))$ for all $v$. Let $v \sim w$, so $\pi(v) = \pi(w)$. Then:
\[
P(v) = f(\pi(v)) = f(\pi(w)) = P(w)
\]
\end{proof}

\begin{remark}[Information-theoretic nature]
The barrier is \emph{informational}, not computational. Given unlimited time, memory, and computational power, an attribute-only observer still cannot distinguish $v$ from $w$ when $\pi(v) = \pi(w)$. The constraint is on the evidence itself.
\end{remark}

\begin{remark}[Role in the paper]
Theorem~\ref{thm:information-barrier} is the base invariance statement. The technical contribution is the downstream structure built on top of it: the ambiguity-based converse (Theorem~\ref{thm:converse}), the Pareto characterization (Theorem~\ref{thm:lwd-optimal}), and the matroid/equicardinality results (Section~\ref{sec:matroid}).
\end{remark}

\begin{corollary}[Class identity is not attribute-computable]\label{cor:provenance-barrier}
Let $C: \mathcal{V} \to \{1,\ldots,k\}$ be the class assignment function. If there exist values $v, w$ with $\pi(v) = \pi(w)$ but $C(v) \neq C(w)$, then class identity is not attribute-computable.
\end{corollary}
\leanmeta{ \LH{ACS9}}

\begin{proof}
Direct application of Theorem~\ref{thm:information-barrier} to $P = C$.
\end{proof}

\subsection{The Positive Result: Nominal Tagging}

We now show that augmenting attribute observations with a single primitive, nominal-tag access, achieves constant witness cost.

\begin{definition}[Nominal-tag access]
A \emph{nominal tag} is a value $\tau(v) \in \mathcal{T}$ associated with each $v \in \mathcal{V}$, representing the class identity of $v$. The \emph{nominal-tag access} operation returns $\tau(v)$ in $O(1)$ time.
\end{definition}

\begin{definition}[Primitive query set]
The extended primitive query set is $\Phi_{\mathcal{I}}^+ = \Phi_{\mathcal{I}} \cup \{\tau\}$, where $\tau$ denotes nominal-tag access.
\end{definition}

\begin{definition}[Witness cost]
Let $W(P)$ denote the minimum number of primitive queries from $\Phi_{\mathcal{I}}^+$ required to compute property $P$. We distinguish two tasks:
\begin{itemize}
\item $W_{\text{id}}$: Cost to identify the class of a single entity.
\item $W_{\text{eq}}$: Cost to determine if two entities have the same class.
\end{itemize}
Unless specified, $W$ refers to $W_{\text{eq}}$.
\end{definition}

\begin{theorem}[Constant witness for class identity]\label{thm:constant-witness}
Under nominal-tag access, class identity checking has constant witness cost:
\[
W(\text{class-identity}) = O(1)
\]
Specifically, the witness procedure is: return $\tau(v_1) = \tau(v_2)$.
\end{theorem}
\leanmeta{ \LH{ACS7}, \LH{FORC2}, \LH{DOM1}}

\begin{proof}
The procedure makes exactly 2 primitive queries (one $\tau$ access per value) and one comparison. This is $O(1)$ regardless of the number of available attributes $|\mathcal{I}|$.
\end{proof}

\begin{theorem}[Attribute-only lower bound]\label{thm:interface-lower-bound}
For attribute-only observers, class identity checking requires:
\[
W(\text{class-identity}) = \Omega(d)
\]
in the worst case, where $d$ is the distinguishing dimension (Definition~\ref{def:distinguishing-dimension}).
\end{theorem}
\leanmeta{ \LH{ACS1}}

\begin{proof}
Assume a zero-error attribute-only procedure halts after fewer than $d$ queries on every execution path. Fix any execution path and let $Q \subseteq \mathcal{I}$ be the set of queried attributes on that path, so $|Q|<d$. Since $d$ is the cardinality of every minimal distinguishing set, no set of size $<d$ is distinguishing; hence there exist values $v,w$ from different classes with identical answers on all attributes in $Q$.

An adversary can answer the procedure's queries consistently with both $v$ and $w$ along this path. Therefore the resulting transcript (and output) is identical on $v$ and $w$, contradicting zero-error class identification. So some execution path must use at least $d$ queries, giving worst-case cost $\Omega(d)$.
\end{proof}

\subsection{Main Contributions}

This paper establishes the following results:

\begin{enumerate}
\item \textbf{Information Barrier Theorem} (Theorem~\ref{thm:information-barrier}): Attribute-only observers cannot compute any property that varies within $\sim$-equivalence classes. This is an observational impossibility, not a computational limitation, and it is the source of the downstream rate-query tradeoffs.

\item \textbf{Constant-Witness Theorem} (Theorem~\ref{thm:constant-witness}): Nominal-tag access achieves constant witness cost for class identity. Attribute-only observers have matching lower bound $\Omega(d)$ (Theorem~\ref{thm:interface-lower-bound}), where $d$ is the distinguishing dimension (Definition~\ref{def:distinguishing-dimension}).

\item \textbf{Complexity Separation} (Section~\ref{sec:complexity}): We establish O(1) vs O(k) vs $\Omega(d)$ complexity bounds for error localization under different observation regimes (where $d$ is the distinguishing dimension).

\item \textbf{Matroid Structure} (Section~\ref{sec:matroid}): Minimal distinguishing query sets form the bases of a matroid. All such sets have equal cardinality, so the distinguishing dimension is an invariant of the observation model rather than an artifact of a particular proof or query order.

\item \textbf{$(L, W, D)$ Optimality} (Section~\ref{sec:lwd}): We characterize the zero-error converse via collision multiplicity $A_\pi$ and prove uniqueness of the nominal point in the maximal-barrier regime ($A_\pi=k$), turning an informal heuristic into an explicit resource law.

\item \textbf{Finite-Block Confusability Law} (Section~\ref{sec:framework}): We prove exact block thresholds for zero-error feasibility in the confusability formulation: at blocklength $t$, feasibility is equivalent to an alphabet threshold $A_\pi^t$, with exact linear log-bit scaling and per-coordinate stabilization.

\item \textbf{Open-World Robustness + Decidability Boundary}: We formalize that barrier-freedom is not extension-stable in open worlds (any finite barrier-free world can be extended to force a collision), and we prove a Rice-style impossibility theorem: no computable certifier can decide barrier existence/freedom for arbitrary generators.

\item \textbf{Machine-Checked Proofs}: All results formalized in Lean 4 (\LeanTotalLines\ lines, \LeanTotalTheorems\ theorem/lemma statements, \LeanTotalSorry\ \texttt{sorry} placeholders).
\end{enumerate}
\leanmeta{ \LH{ROB1}, \LH{ROB2}, \LH{UND1}, \LH{UND2} }

Without the equicardinality result behind the matroid structure, the lower bound would depend on which distinguishing set or query strategy one chose. The matroid theorem is what makes $d$ strategy-independent.

\subsection{Related Work and Positioning}

This paper does not develop asymptotic coding theorems or claim new zero-error capacity results. Its contribution is to make explicit the bit/query/error tradeoffs governing zero-error classification under constrained observations.

\textbf{Identification via channels.} Our work is adjacent to the identification paradigm introduced by Ahlswede and Dueck \cite{ahlswede1989identification, ahlswede1989identification2}. In their framework, a decoder need not reconstruct a message but only answer ``is the message $m$?'' for a given hypothesis. Our setting is different in three ways: (1) we consider one-shot zero-error class identification rather than asymptotic vanishing-error identification, (2) queries are adaptive rather than block codes, and (3) we allow auxiliary tagging (rate $L$) to reduce query cost. The connection is motivational and terminological rather than a direct reuse of the original coding theorem.

\textbf{Rate-distortion theory.} The $(L, W, D)$ framework is analogous to rate-distortion theory \cite{shannon1959coding, berger1971rate}: the ``distortion'' $D$ is semantic (class misidentification), and there is a second resource $W$ (query cost) alongside rate $L$. Classical rate-distortion asks for the minimum rate needed to achieve distortion $D$; here the main question is which combinations of tag bits and queries permit $D=0$. Theorem~\ref{thm:lwd-optimal} gives the resulting converse in terms of collision multiplicity and identifies the unique nominal point in the maximal-barrier regime.

\textbf{Rate-distortion-perception tradeoffs.} Blau and Michaeli \cite{blau2019rethinking} extended rate-distortion theory by adding a perception constraint, creating a three-way tradeoff. Our query cost $W$ plays an analogous role: it measures the interactive cost of achieving low distortion rather than a distributional constraint. We cite this as a structural analogy, not as a claim that the present model inherits the full geometry of that setting.

\textbf{Zero-error information theory.} The matroid structure (Section~\ref{sec:matroid}) places the model near zero-error information theory. K\"orner \cite{korner1973coding} and Witsenhausen \cite{witsenhausen1976zero} studied zero-error source coding where confusable symbols must be distinguished. Here the direct theorem is simpler: when $L = 0$, the distinguishing dimension (Definition~\ref{def:distinguishing-dimension}) is the minimum number of binary queries needed for zero-error class separation. In addition, the paper proves exact finite-block confusability laws: block feasibility threshold $A_\pi^k$, exact linear log-bit scaling, and per-coordinate stabilization in the worst-case regime. These are formalized in Lean and strengthen the operational converse while remaining distinct from full distribution-optimized graph-entropy/capacity theorems.

\textbf{Query complexity and communication complexity.} The $\Omega(d)$ lower bound for attribute-only identification relates to decision tree complexity \cite{buhrman2002complexity} and interactive communication \cite{orlitsky1991worst}. The key distinction is that our queries are constrained to a fixed observable family $\mathcal{I}$, not arbitrary predicates.

\textbf{Compression in classification systems.} The same ambiguity bound appears across databases, knowledge graphs, taxonomy, and software-runtime typing systems: for zero-error identification, ambiguity induces a minimum metadata requirement $L \ge \log_2 A_\pi$ (Theorem~\ref{thm:converse}), with maximal-barrier specialization $L \ge \log_2 k$.

\textbf{Runtime/schema instantiation (secondary).} In nominal-vs-structural interface settings \cite{Cardelli1985, cook1990inheritance}, the model yields a concrete cost statement: under attribute collisions, purely profile-based identification has worst-case $\Omega(d)$ witness cost, while explicit tags achieve $O(1)$ identification using $O(\log A_\pi)$ bits. This is one instantiation of the general bounds.

\subsection{Paper Organization}

Section~\ref{sec:framework} formalizes the compression framework and defines the $(L, W, D)$ tradeoff. Section~\ref{sec:complexity} establishes complexity bounds for error localization. Section~\ref{sec:matroid} proves the matroid structure of distinguishing query families. Section~\ref{sec:witness} analyzes witness cost in detail. Section~\ref{sec:lwd} proves the ambiguity-based converse and Pareto characterization. Section~\ref{sec:conclusion} concludes. Extended instantiations and the Lean formalization are provided in the supplementary material.
